\chapter{Pulmonary Artery (Swan-Ganz) Catheterization}

\section*{Overview}
Pulmonary artery catheterization is the placement of a balloon-tipped catheter through the right heart into the pulmonary artery to measure hemodynamic pressures and cardiac output. It provides detailed right and left heart hemodynamic information at the bedside.

\section*{Alternative Names}
Swan-Ganz catheterization, right heart catheterization, pulmonary artery catheter, hemodynamic monitoring

\section*{Principle}
The catheter, floated through the right atrium and ventricle with an inflated balloon, wedges in a distal pulmonary artery, allowing measurement of pulmonary capillary wedge pressure, a surrogate for left atrial and left ventricular end-diastolic pressure. Pressure waveforms, mixed venous oxygen saturation, and thermodilution cardiac output are used to characterize shock, guide fluid and vasopressor therapy, and assess ventricular function.

\section*{History}
The flow-directed balloon catheter was introduced by Swan and Ganz in 1970 and revolutionized bedside hemodynamic monitoring in intensive care. Its use declined after trials questioned benefit in unselected patients, but it remains valuable in selected complex cases.

\section*{Why It Is Done}
Ordered when bedside assessment of hemodynamics is insufficient, including undifferentiated shock, severe heart failure with unclear volume status, pulmonary hypertension, and high-risk cardiac surgery. It helps distinguish cardiogenic, distributive, and hypovolemic shock and guides titration of fluids, inotropes, and vasopressors.

\section*{Is This a Primary or Secondary Test or Procedure}
A secondary, specialized diagnostic and monitoring procedure reserved for patients in whom standard clinical and noninvasive assessment cannot establish hemodynamics.

\section*{How It Is Performed}
A central venous sheath is placed, and the catheter is advanced to the right atrium where the balloon is inflated and floated through the right ventricle into the pulmonary artery under continuous waveform and pressure monitoring. The balloon is deflated and the catheter left in place, connected to transducers and a cardiac output computer. Waveforms confirm sequential positioning in the right atrium, right ventricle, and pulmonary artery, then a wedge tracing on inflation.

\section*{Preparation}
Informed consent, coagulation assessment, and continuous electrocardiographic and blood pressure monitoring. Sterile technique is used for central venous access, and the patient is positioned supine.

\section*{Contraindications and Precautions}
Relative contraindications include severe coagulopathy, left bundle branch block (risk of complete heart block), recent cardiac surgery with tricuspid valve or right heart pathology, and known right-sided endocarditis. Precautions include avoiding prolonged balloon inflation, monitoring for arrhythmias and catheter coiling, and using ultrasound for venous access.

\section*{Risks}
Arrhythmias, pulmonary artery rupture or infarction, pneumothorax and bleeding from venous access, infection, catheter knotting, and pulmonary thromboembolism.

\section*{Aftercare and Recovery}
Hemodynamic parameters are recorded and used to guide therapy, with the catheter generally removed within 24--72 hours. The insertion site is observed for bleeding and infection, and monitoring is continued until the catheter is withdrawn.

\section*{Outcome and Follow-up}
Cardiogenic shock shows low cardiac output with elevated wedge pressure and filling pressures, while distributive shock shows low or normal wedge pressure with high output. A low wedge pressure suggests hypovolemia warranting fluid resuscitation, whereas a high wedge pressure in a hypotensive patient supports diuresis or inotropic support.

\section*{Expected Results}
Normal values include cardiac index of 2.5--4.0 L/min/m$^2$, pulmonary artery systolic pressure of 20--30 mm Hg, and mean pulmonary capillary wedge pressure of 6--12 mm Hg.

\section*{Follow-up}
Repeat cardiac output and pressure measurements during therapy guide further management. Echocardiography and other imaging may be performed to corroborate findings, and the catheter is removed once hemodynamics stabilize or noninvasive monitoring suffices.

\section*{When to Seek Medical Care}
Follow the procedure team's specific instructions. Seek urgent medical assessment for severe or worsening pain, uncontrolled bleeding, fever or signs of infection, trouble breathing, chest pain, fainting, new weakness or confusion, inability to urinate, or any rapidly worsening symptom. The appropriate warning signs depend on the procedure and the patient's medical condition.

\section*{References}
\begin{enumerate}
  \item MedlinePlus. Pulmonary artery catheterization. U.S. National Library of Medicine; 2025.
  \item Current Medical Diagnosis and Treatment. McGraw-Hill; 2025.
\end{enumerate}

\clearpage
